% ==============================================================================
% Carbon Pricing in Reverse: Evidence from Virginia's Exit and Re-Entry into RGGI
% ------------------------------------------------------------------------------
% Engine: pdfLaTeX.  Compile from this directory (Paper/) with:
%   TEXINPUTS=../Preambles:$TEXINPUTS pdflatex -interaction=nonstopmode rggi_carbon_pricing_reverse.tex
%   BIBINPUTS=..:$BIBINPUTS bibtex rggi_carbon_pricing_reverse
%   TEXINPUTS=../Preambles:$TEXINPUTS pdflatex -interaction=nonstopmode rggi_carbon_pricing_reverse.tex
%   TEXINPUTS=../Preambles:$TEXINPUTS pdflatex -interaction=nonstopmode rggi_carbon_pricing_reverse.tex
%
% Structure rule: ALL numbers, tables, and figures are produced by scripts/R/*
% and pulled in via \input{../Tables/...} and \includegraphics{...}.
% Do NOT hand-type estimates. See quality_reports/plans/2026-05-30_rggi-paper.md.
% ==============================================================================
\documentclass[12pt]{article}
% ==============================================================================
% header.tex -- shared preamble for the RGGI paper (article class, pdfLaTeX)
% ==============================================================================
% Engine: pdfLaTeX (XeLaTeX not required; avoid fontspec-only features).
% Included by Paper/*.tex via:  \documentclass[12pt]{article}% ==============================================================================
% header.tex -- shared preamble for the RGGI paper (article class, pdfLaTeX)
% ==============================================================================
% Engine: pdfLaTeX (XeLaTeX not required; avoid fontspec-only features).
% Included by Paper/*.tex via:  \documentclass[12pt]{article}% ==============================================================================
% header.tex -- shared preamble for the RGGI paper (article class, pdfLaTeX)
% ==============================================================================
% Engine: pdfLaTeX (XeLaTeX not required; avoid fontspec-only features).
% Included by Paper/*.tex via:  \documentclass[12pt]{article}\input{header}
% Compile from Paper/ with: TEXINPUTS=../Preambles:$TEXINPUTS pdflatex ...
% ==============================================================================

% --- Encoding & fonts (pdfLaTeX) ---------------------------------------------
\usepackage[T1]{fontenc}
\usepackage[utf8]{inputenc}
\usepackage{lmodern}
\usepackage{microtype}

% --- Page geometry & spacing -------------------------------------------------
\usepackage[margin=1in]{geometry}
\usepackage{setspace}
\usepackage{parskip}        % author may switch to \onehalfspacing for review

% --- Math --------------------------------------------------------------------
\usepackage{amsmath,amssymb,amsthm,mathtools}
\usepackage{bm}

% --- Tables & figures --------------------------------------------------------
\usepackage{graphicx}
\usepackage{booktabs}
\usepackage{threeparttable}  % table notes under regression tables
\usepackage{dcolumn}         % decimal-aligned columns
\usepackage{siunitx}
\usepackage{multirow}
\usepackage{array}
\usepackage{float}
\usepackage{caption}
\usepackage{subcaption}
\usepackage{rotating}        % sidewaystable for wide results
\usepackage{longtable}
\graphicspath{{../Figures/}}

% --- Bibliography (natbib, econ author-year) ---------------------------------
\usepackage{natbib}
\bibliographystyle{plainnat} % swap to aer.bst / a journal style at submission
\setcitestyle{authoryear,round,aysep={},yysep={;}}

% --- Cross-refs & links ------------------------------------------------------
\usepackage[colorlinks=true,allcolors=blue!55!black,breaklinks=true]{hyperref}
\usepackage{cleveref}

% --- Theorem-like environments ------------------------------------------------
\theoremstyle{plain}
\newtheorem{prediction}{Prediction}
\theoremstyle{definition}
\newtheorem{assumption}{Assumption}

% --- Drafting helpers (strip before submission) ------------------------------
\usepackage{xcolor}
\newcommand{\TODO}[1]{\textcolor{red}{\textbf{[TODO: #1]}}}
\newcommand{\PENDING}[1]{\textcolor{orange}{\textbf{[PENDING DATA: #1]}}}

% ==============================================================================
% Project notation macros
% ==============================================================================
% Treatment / estimands
\newcommand{\VA}{\text{VA}}
\newcommand{\treat}{D}                 % treatment indicator
\newcommand{\effect}{\tau}             % treatment effect / ATT
\newcommand{\effexit}{\hat{\effect}^{\text{exit}}}
\newcommand{\effentry}{\hat{\effect}^{\text{entry}}}

% Carbon-adder partial-equilibrium model (Sec. 4)
\newcommand{\pretail}{p_{r}}           % retail price
\newcommand{\cgen}{c_{g}}              % generation cost
\newcommand{\ctd}{c_{T}}               % transmission & distribution cost
\newcommand{\pallow}{\pi_{\text{allow}}} % allowance price ($/ton)
\newcommand{\erate}{e}                 % emissions rate of marginal generator
\newcommand{\markup}{\mu}              % regulated markup
\newcommand{\co}{\mathrm{CO}_2}

% Units
\newcommand{\dollarMWh}{\$/\mathrm{MWh}}
\newcommand{\tonsMWh}{\mathrm{tons}/\mathrm{MWh}}

% Abbreviations
\newcommand{\rggi}{RGGI}
\newcommand{\pjm}{PJM}
\newcommand{\scm}{SCM}
\newcommand{\did}{DiD}
\newcommand{\sdid}{SDiD}

% Compile from Paper/ with: TEXINPUTS=../Preambles:$TEXINPUTS pdflatex ...
% ==============================================================================

% --- Encoding & fonts (pdfLaTeX) ---------------------------------------------
\usepackage[T1]{fontenc}
\usepackage[utf8]{inputenc}
\usepackage{lmodern}
\usepackage{microtype}

% --- Page geometry & spacing -------------------------------------------------
\usepackage[margin=1in]{geometry}
\usepackage{setspace}
\usepackage{parskip}        % author may switch to \onehalfspacing for review

% --- Math --------------------------------------------------------------------
\usepackage{amsmath,amssymb,amsthm,mathtools}
\usepackage{bm}

% --- Tables & figures --------------------------------------------------------
\usepackage{graphicx}
\usepackage{booktabs}
\usepackage{threeparttable}  % table notes under regression tables
\usepackage{dcolumn}         % decimal-aligned columns
\usepackage{siunitx}
\usepackage{multirow}
\usepackage{array}
\usepackage{float}
\usepackage{caption}
\usepackage{subcaption}
\usepackage{rotating}        % sidewaystable for wide results
\usepackage{longtable}
\graphicspath{{../Figures/}}

% --- Bibliography (natbib, econ author-year) ---------------------------------
\usepackage{natbib}
\bibliographystyle{plainnat} % swap to aer.bst / a journal style at submission
\setcitestyle{authoryear,round,aysep={},yysep={;}}

% --- Cross-refs & links ------------------------------------------------------
\usepackage[colorlinks=true,allcolors=blue!55!black,breaklinks=true]{hyperref}
\usepackage{cleveref}

% --- Theorem-like environments ------------------------------------------------
\theoremstyle{plain}
\newtheorem{prediction}{Prediction}
\theoremstyle{definition}
\newtheorem{assumption}{Assumption}

% --- Drafting helpers (strip before submission) ------------------------------
\usepackage{xcolor}
\newcommand{\TODO}[1]{\textcolor{red}{\textbf{[TODO: #1]}}}
\newcommand{\PENDING}[1]{\textcolor{orange}{\textbf{[PENDING DATA: #1]}}}

% ==============================================================================
% Project notation macros
% ==============================================================================
% Treatment / estimands
\newcommand{\VA}{\text{VA}}
\newcommand{\treat}{D}                 % treatment indicator
\newcommand{\effect}{\tau}             % treatment effect / ATT
\newcommand{\effexit}{\hat{\effect}^{\text{exit}}}
\newcommand{\effentry}{\hat{\effect}^{\text{entry}}}

% Carbon-adder partial-equilibrium model (Sec. 4)
\newcommand{\pretail}{p_{r}}           % retail price
\newcommand{\cgen}{c_{g}}              % generation cost
\newcommand{\ctd}{c_{T}}               % transmission & distribution cost
\newcommand{\pallow}{\pi_{\text{allow}}} % allowance price ($/ton)
\newcommand{\erate}{e}                 % emissions rate of marginal generator
\newcommand{\markup}{\mu}              % regulated markup
\newcommand{\co}{\mathrm{CO}_2}

% Units
\newcommand{\dollarMWh}{\$/\mathrm{MWh}}
\newcommand{\tonsMWh}{\mathrm{tons}/\mathrm{MWh}}

% Abbreviations
\newcommand{\rggi}{RGGI}
\newcommand{\pjm}{PJM}
\newcommand{\scm}{SCM}
\newcommand{\did}{DiD}
\newcommand{\sdid}{SDiD}

% Compile from Paper/ with: TEXINPUTS=../Preambles:$TEXINPUTS pdflatex ...
% ==============================================================================

% --- Encoding & fonts (pdfLaTeX) ---------------------------------------------
\usepackage[T1]{fontenc}
\usepackage[utf8]{inputenc}
\usepackage{lmodern}
\usepackage{microtype}

% --- Page geometry & spacing -------------------------------------------------
\usepackage[margin=1in]{geometry}
\usepackage{setspace}
\usepackage{parskip}        % author may switch to \onehalfspacing for review

% --- Math --------------------------------------------------------------------
\usepackage{amsmath,amssymb,amsthm,mathtools}
\usepackage{bm}

% --- Tables & figures --------------------------------------------------------
\usepackage{graphicx}
\usepackage{booktabs}
\usepackage{threeparttable}  % table notes under regression tables
\usepackage{dcolumn}         % decimal-aligned columns
\usepackage{siunitx}
\usepackage{multirow}
\usepackage{array}
\usepackage{float}
\usepackage{caption}
\usepackage{subcaption}
\usepackage{rotating}        % sidewaystable for wide results
\usepackage{longtable}
\graphicspath{{../Figures/}}

% --- Bibliography (natbib, econ author-year) ---------------------------------
\usepackage{natbib}
\bibliographystyle{plainnat} % swap to aer.bst / a journal style at submission
\setcitestyle{authoryear,round,aysep={},yysep={;}}

% --- Cross-refs & links ------------------------------------------------------
\usepackage[colorlinks=true,allcolors=blue!55!black,breaklinks=true]{hyperref}
\usepackage{cleveref}

% --- Theorem-like environments ------------------------------------------------
\theoremstyle{plain}
\newtheorem{prediction}{Prediction}
\theoremstyle{definition}
\newtheorem{assumption}{Assumption}

% --- Drafting helpers (strip before submission) ------------------------------
\usepackage{xcolor}
\newcommand{\TODO}[1]{\textcolor{red}{\textbf{[TODO: #1]}}}
\newcommand{\PENDING}[1]{\textcolor{orange}{\textbf{[PENDING DATA: #1]}}}

% ==============================================================================
% Project notation macros
% ==============================================================================
% Treatment / estimands
\newcommand{\VA}{\text{VA}}
\newcommand{\treat}{D}                 % treatment indicator
\newcommand{\effect}{\tau}             % treatment effect / ATT
\newcommand{\effexit}{\hat{\effect}^{\text{exit}}}
\newcommand{\effentry}{\hat{\effect}^{\text{entry}}}

% Carbon-adder partial-equilibrium model (Sec. 4)
\newcommand{\pretail}{p_{r}}           % retail price
\newcommand{\cgen}{c_{g}}              % generation cost
\newcommand{\ctd}{c_{T}}               % transmission & distribution cost
\newcommand{\pallow}{\pi_{\text{allow}}} % allowance price ($/ton)
\newcommand{\erate}{e}                 % emissions rate of marginal generator
\newcommand{\markup}{\mu}              % regulated markup
\newcommand{\co}{\mathrm{CO}_2}

% Units
\newcommand{\dollarMWh}{\$/\mathrm{MWh}}
\newcommand{\tonsMWh}{\mathrm{tons}/\mathrm{MWh}}

% Abbreviations
\newcommand{\rggi}{RGGI}
\newcommand{\pjm}{PJM}
\newcommand{\scm}{SCM}
\newcommand{\did}{DiD}
\newcommand{\sdid}{SDiD}


\title{Carbon Pricing in Reverse:\\
       Evidence from Virginia's Exit and Re-Entry into RGGI}

\author{%
  \TODO{Author Name}\\
  \TODO{Affiliation}\\
  \texttt{\TODO{email}}%
}

\date{\today \\[0.5em] \footnotesize Preliminary draft --- please do not cite or circulate.}

\begin{document}
\maketitle

\begin{abstract}
\noindent
\PENDING{Abstract.} Virginia is the only U.S. state to enter, exit, and re-enter a
carbon cap-and-trade program (the Regional Greenhouse Gas Initiative, RGGI) within a
five-year span. Because the exit was driven by a change in gubernatorial administration
rather than by economic conditions, the timing of treatment is plausibly exogenous. We
use synthetic control and event-study difference-in-differences methods on public data to
estimate the causal effect of Virginia's RGGI participation on retail electricity prices,
wholesale prices, and \co{} emissions, and we test whether the effects of \emph{exiting}
a carbon market are symmetric to the effects of \emph{entering} it. \PENDING{Headline
estimates for the entry (Jan.\ 2021) and exit (Dec.\ 2023) events; the re-entry event
(Jul.\ 2026) is analyzed prospectively.}

\medskip
\noindent\textbf{Keywords:} carbon pricing; cap-and-trade; RGGI; synthetic control;
electricity markets; emissions leakage; policy reversal.\\
\noindent\textbf{JEL Codes:} Q48, Q54, Q58, L94, H23.
\end{abstract}

\clearpage
\tableofcontents
\clearpage

% ==============================================================================
\section{Introduction}\label{sec:intro}
% --- Sec 1.1 Motivation; 1.2 Research Questions; 1.3 Contribution ---
\TODO{Phase 5 prose.} Carbon cap-and-trade is the dominant market-based climate policy
instrument in the United States, yet credible causal evidence on its effects is scarce
because clean counterfactuals are rare. Virginia provides a unique natural experiment.

\paragraph{Research questions.} \TODO{prose around the four questions:}
(i) the causal effect of Virginia's RGGI exit on retail prices, wholesale prices, and
\co{} emissions; (ii) symmetry of exit vs.\ entry effects; (iii) the extent of within-PJM
emissions leakage; (iv) whether re-entry replicates the entry effects.

\paragraph{Contribution.} \TODO{prose:} first empirical study of a carbon-market
\emph{exit}; a test for asymmetry/hysteresis exploiting a rare reversal; within-RTO (PJM)
leakage estimates under sub-national pricing; a paired entry/exit/re-entry design on a
single treated unit. Early illustrative evidence appears in \citet{Murray2015_rggi_attribution}.

% ==============================================================================
\section{Institutional Background}\label{sec:background}

\subsection{RGGI Overview}\label{sec:rggi-overview}
\TODO{Phase 5 prose:} 2009 launch, program reviews, current membership; quarterly
allowance auctions, compliance obligations, Cost Containment Reserve (CCR), Emissions
Containment Reserve (ECR). Third Program Review changes effective 2027 (tighter cap
trajectory, higher minimum reserve price, elimination of offsets).

\subsection{Virginia's RGGI Timeline}\label{sec:va-timeline}
\begin{table}[H]
\centering
\caption{Virginia's RGGI participation timeline.}
\label{tab:va-timeline}
\begin{threeparttable}
\begin{tabular}{@{}p{2.6cm}p{6.2cm}p{4.6cm}@{}}
\toprule
Date & Event & Political context \\
\midrule
Apr.\ 2020 & Clean Economy Act \& Community Flood Preparedness Act signed & Democratic trifecta \\
Jan.\ 2021 & Virginia begins RGGI participation & First compliance period \\
Jan.\ 2022 & Youngkin inaugurated; signals intent to withdraw & Republican governor, split legislature \\
Jun.\ 2023 & Air Pollution Control Board votes to repeal RGGI regulation & Executive-driven withdrawal \\
Dec.\ 2023 & Last Virginia participation in RGGI auction & Formal exit \\
Jan.\ 2024--Jun.\ 2026 & Virginia outside RGGI & Non-participation period \\
2025 & Spanberger (D) elected; legislation to rejoin signed & Democratic trifecta restored \\
Jul.\ 2026 & Virginia re-enters RGGI & Second participation period \\
\bottomrule
\end{tabular}
\begin{tablenotes}\footnotesize
\item \textit{Sources:} \TODO{Phase 5 -- cite Clean Economy Act; APCB repeal; RGGI auction records.}
\end{tablenotes}
\end{threeparttable}
\end{table}

\subsection{The Dominion Energy Context}\label{sec:dominion}
\TODO{Phase 5 prose, public sources only:} Dominion-dominated regulated market; RGGI costs
passed through via a rider (\TODO{rider \$/month from public SCC dockets}); rider removed on
exit, to be reinstated on re-entry. Concurrent factors: residential rate increases, data
center load growth, IRA-driven clean-energy investment. \emph{No Dominion internal data are
used; rider figures are drawn from public SCC filings.}

% ==============================================================================
\section{Literature Review}\label{sec:lit}
\TODO{Phase 5 prose, organized in four strands:}
\begin{itemize}
  \item \textbf{Empirical effects of carbon pricing:} \citet{Murray2015_rggi_attribution}; \TODO{Fell \& Maniloff 2018; Yan 2021; Song \& Hochman 2025.}
  \item \textbf{Pass-through \& incidence:} \TODO{Fabra \& Reguant 2014; Hintermann 2016; Fowlie, Reguant \& Ryan 2016; Bushnell et al.\ 2013.}
  \item \textbf{Policy reversal \& asymmetry:} \TODO{Greenstone 2002; Coglianese et al.\ 2020; Walker 2013.}
  \item \textbf{Methodology:} \citet{Abadie2010_synth}; \TODO{Abadie et al.\ 2015; Abadie 2021; }\citet{Arkhangelsky2021_sdid}\TODO{; Cattaneo et al.\ 2021; Callaway \& Sant'Anna 2021; Cameron et al.\ 2008.}
\end{itemize}

% ==============================================================================
\section{Theoretical Framework}\label{sec:theory}

\subsection{A Carbon Adder in a Regulated Market}\label{sec:model}
We model the retail price of electricity as
\begin{equation}\label{eq:retail}
  \pretail \;=\; \cgen \;+\; \ctd \;+\; \pallow \cdot \erate \;+\; \markup,
\end{equation}
where $\cgen$ is generation cost (fuel, O\&M), $\ctd$ is transmission and distribution
cost, $\pallow$ is the allowance price ($\dollarMWh$ once scaled by the emissions rate
$\erate$ in $\tonsMWh$ of the marginal generator), and $\markup$ is the regulated markup.
Removing the carbon adder ($\pallow \to 0$) reduces $\pretail$ mechanically through the
rider; the wholesale-price effect depends on whether RGGI-state generators are marginal in
PJM dispatch.

\subsection{Predictions}\label{sec:predictions}
\begin{prediction}[Retail prices]\label{pr:retail}
Retail prices fall by approximately the rider amount on exit and rise by approximately the
new rider on re-entry.
\end{prediction}
\begin{prediction}[Wholesale LMPs]\label{pr:wholesale}
Ambiguous: Virginia generators offer lower prices, but if inframarginal the LMP may not move.
\end{prediction}
\begin{prediction}[Virginia emissions]\label{pr:va-emissions}
Rise if the adder was constraining dispatch of Virginia fossil generators.
\end{prediction}
\begin{prediction}[System-wide emissions]\label{pr:system-emissions}
Ambiguous: depends on whether Virginia generation displaces cleaner or dirtier out-of-state generation (leakage).
\end{prediction}
\begin{prediction}[Asymmetry/hysteresis]\label{pr:asymmetry}
Entry effects may exceed exit effects if 2021--2023 clean investments are irreversible,
efficiency gains persist, or adjustments have ratchet properties.
\end{prediction}

% ==============================================================================
\section{Data}\label{sec:data}
\TODO{Phase 5 prose.} \emph{Public sources only.}
% Placeholder -- regenerated by scripts/R/07_figures_tables.R in Phase 2-3.
% Public data sources only (no Dominion internal data).
\begin{table}[H]
\centering
\caption{Data sources (public). \TODO{Regenerated from config in Phase 2--3.}}
\label{tab:data-sources}
\begin{threeparttable}
\begin{tabular}{@{}p{4.8cm}p{4.2cm}p{2.2cm}p{2.2cm}@{}}
\toprule
Variable & Source & Frequency & Coverage \\
\midrule
Retail electricity prices & EIA Form 861/861M & Monthly/Annual & States, 2018--2026 \\
Wholesale prices (LMPs) & PJM Data Miner 2.0 & Hourly $\to$ monthly & PJM nodes/zones \\
\co{} emissions & EPA CEMS (CAMPD) & Hourly $\to$ monthly & U.S.\ fossil units \\
Generation by fuel & EIA 923 / EPA CEMS & Monthly & U.S.\ generators \\
Allowance auction prices & RGGI Inc.\ reports & Quarterly & 2008--present \\
Natural gas prices & EIA (Henry Hub, Transco Z5) & Daily/Monthly & National/Regional \\
Weather (HDD/CDD) & NOAA ISD/GHCN & Daily & By state \\
Economic controls & BEA / Census & Annual/Quarterly & States \\
Dominion rider & Public SCC dockets & Per filing & Virginia \\
\bottomrule
\end{tabular}
\begin{tablenotes}\footnotesize
\item \textit{Notes:} Public sources only; no Dominion internal data. \TODO{Add access dates/vintages in Phase 2.}
\end{tablenotes}
\end{threeparttable}
\end{table}
   % generated by scripts/R/07_figures_tables.R
% (placeholder file committed so the skeleton compiles; replaced in Phase 2-3)

\subsection{Sample Construction}\label{sec:sample}
\TODO{Phase 5 prose:} unit = state (PJM zone for wholesale); time = month (2018--2026);
treated = Virginia. Control groups: non-RGGI PJM states (WV, OH, PA, KY, IN, NC, TN) and
continuous-RGGI states (MD, DE, NJ, CT, MA, ME, NH, VT, RI, NY).

% ==============================================================================
\section{Empirical Strategy}\label{sec:strategy}

\subsection{Synthetic Control (Primary)}\label{sec:scm}
\TODO{Phase 5 prose:} construct synthetic Virginia matching pre-treatment outcomes and
covariates, separately for the entry (pre: Jan.\ 2018--Dec.\ 2020) and exit (pre: Jan.\
2021--Sep.\ 2023) windows. Inference via placebo/permutation tests and conformal prediction
intervals.

\subsection{Event-Study Difference-in-Differences}\label{sec:did}
\begin{equation}\label{eq:eventstudy}
  Y_{st} \;=\; \alpha_s + \gamma_t
  + \sum_{\tau \neq -1} \beta_\tau \,\mathbf{1}[s=\VA]\,\mathbf{1}[t-t^\ast=\tau]
  + X_{st}'\delta + \varepsilon_{st},
\end{equation}
with state fixed effects $\alpha_s$, month fixed effects $\gamma_t$, time-varying controls
$X_{st}$ (gas prices, HDD, CDD), and dynamic effects $\beta_\tau$ relative to $\tau=-1$.
Standard errors clustered at the state level; wild cluster bootstrap given the small number
of clusters. Estimated separately for the exit ($t^\ast=$ Dec.\ 2023/Jan.\ 2024), entry
($t^\ast=$ Jan.\ 2021), and prospectively re-entry ($t^\ast=$ Jul.\ 2026) events.

\subsection{Synthetic Difference-in-Differences}\label{sec:sdid}
\TODO{Phase 5 prose:} \citet{Arkhangelsky2021_sdid} as a bridge between SCM and DiD.

\subsection{Extensions}\label{sec:extensions}
\TODO{Phase 5 prose:} wholesale PJM-zone analysis (DOM zone treated); emissions leakage
(VA-plant vs.\ PJM-wide emissions); formal asymmetry test $|\effexit| = |\effentry|$;
distributional analysis if public microdata permit.

% ==============================================================================
\section{Threats to Identification and Robustness}\label{sec:threats}
\begin{table}[H]
\centering
\caption{Threats and robustness checks.}
\label{tab:threats}
\begin{threeparttable}
\begin{tabular}{@{}p{5.4cm}p{8.0cm}@{}}
\toprule
Threat & Test / mitigation \\
\midrule
Anticipation (exit signaled early 2022) & Alternative treatment dates (Jan.\ 2022, Jun.\ 2023, Dec.\ 2023); pre-trend tests. \\
Concurrent policy shocks (IRA, RPS) & Control for renewable capacity additions and federal subsidy proxies. \\
Gas price volatility & Control for regional gas hub prices; gas-adjusted outcomes. \\
Data center load growth & Control for large-customer load; use per-MWh prices. \\
SUTVA violation (PJM spillovers) & Estimate PJM-wide wholesale effects; bound spillovers. \\
Small-$N$ inference & Permutation inference (SCM); wild cluster bootstrap (DiD). \\
Rider mechanicality & Decompose retail effect into rider (mechanical) and non-rider components. \\
Fleet composition changes & Control for capacity additions/retirements; shares vs.\ levels. \\
\bottomrule
\end{tabular}
\end{threeparttable}
\end{table}

% ==============================================================================
\section{Results}\label{sec:results}
\PENDING{All results below are generated by \texttt{scripts/R/} from public data and pulled
in via \texttt{\textbackslash input}/\texttt{\textbackslash includegraphics}. No numbers are
hand-entered.}

\subsection{Exit Event (Dec.\ 2023)}\label{sec:res-exit}
\PENDING{SCM gap plot, event-study coefficients, placebo distribution.}
% \input{../Tables/tab_main_exit.tex}
% \includegraphics[width=\linewidth]{scm_gap_exit_retail.pdf}

\subsection{Entry Event (Jan.\ 2021)}\label{sec:res-entry}
\PENDING{SCM and event-study estimates for entry.}
% \input{../Tables/tab_main_entry.tex}

\subsection{Asymmetry Test}\label{sec:res-asymmetry}
\PENDING{Formal test of $|\effexit| = |\effentry|$.}

\subsection{Emissions Leakage}\label{sec:res-leakage}
\PENDING{VA-plant vs.\ PJM-wide emissions; implied leakage rate.}

% ==============================================================================
\section{Expected Results and Interpretation}\label{sec:interpretation}
\TODO{Phase 5 prose around three pre-registered scenarios:}
\begin{description}
  \item[Scenario A --- Symmetric effects.] Carbon-pricing effects are reversible; policy
    uncertainty may undermine long-run investment incentives.
  \item[Scenario B --- Asymmetric (entry $>$ exit).] Irreversible clean investments create a
    ratchet; even temporary participation yields lasting benefits.
  \item[Scenario C --- Leakage dominates.] Sub-national design reshuffles generation across
    state lines without net reductions.
\end{description}

\paragraph{Re-entry (prospective).} \PENDING{The Jul.\ 2026 re-entry is in the future as of
this draft. We pre-register the design (treatment date, donor pool, estimands, inference) and
will report estimates once post-treatment data are available.}

% ==============================================================================
\section{Conclusion}\label{sec:conclusion}
\TODO{Phase 5 prose.}

% ==============================================================================
\clearpage
% Interim: canonical Bibliography_base.bib is hook-protected; using the parallel
% verified bib in master_supporting_docs/ until that is unprotected (then switch back).
\bibliography{../master_supporting_docs/Bibliography_verified}

% ==============================================================================
\clearpage
\appendix
\section{Data Sources and Public Retrieval Methods}\label{app:data}
\TODO{Phase 3 -- map each variable to its public source and access method.}

\section{Pre-Analysis Plan: Re-Entry (Jul.\ 2026)}\label{app:pap}
\TODO{Phase 4 -- pre-registered estimands, treatment date, donor pool, and inference for the
re-entry event, fixed in advance of observing outcomes.}

\end{document}
